\scanpage{25}
% L01–05
monomorphisme
\[
\mathrm{Aff}(1)_S=G_m\underset{1/2}{\boldsymbol{.}}G_a\xrightarrow{i}G
\]
et une section $t$ de $G_a$, telle que $u=i(t)$. Je dis qu'alors aussi on a
\[
A(u)=2
\]

% L06–12
En effet, on trouve plus \deleted{précisément} \inserted{généralement}
\[
A(i(\lambda:t))=\lambda+\lambda^{-1}
\]
On le voit p. ex. en notant que (pour $\lambda,t$ indét)
\[
f(\lambda,t)\overset{\text{déf}}{=}A(i(\lambda,t)),
\]
doit être un $\in\underline O_S[\lambda,t][\lambda^{-1}]$ satisfaisant les conditions
\[
\begin{aligned}
&\left\{\begin{aligned}
&f(\lambda,0)=\lambda+\lambda^{-1}\\
&\text{plus gén}\quad f(\underbrace{(\unsure{0|\omega},t)(\lambda,0)(0,-t)}_{(\lambda,0)(\lambda^{-1}t-t)=(\lambda,\lambda^{-1}t-t)})=\lambda+\lambda^{-1}
\end{aligned}\right.
\end{aligned}
\]

% L13–17
donc, posant $s=\lambda^{-1}t-t$ i.e. $t=\dfrac{s}{\lambda^{-1}-1}$, on a
\[
f(\lambda,s)=\lambda+\lambda^{-1}
\]
dès que $\lambda^{-1}-1$ inversible -- donc par prolongement, pour tt système $(\lambda,s)$.

% L18–24
Une condition \uline{nécessaire} pour que $u$ soit unipotent est donc que $A(u)=1$. On voit aisément que cette condition est suffisante au dessus d'un corps (utilisant le fait que sur un corps alg. clos, tt élément de $G$ est contenu dans un ss-groupe de Borel). J'ignore ce qui en est dans le cas général.
% archive sheet number: 24
